\documentclass{article}

% IsabelleBlueprint reads the theorem-like environments below directly.
% You do not need to compile this with LaTeX for `isabelle-blueprint` to
% parse it, but it is written to be valid so you can also build a PDF.

\usepackage{amsthm}
\newtheorem{definition}{Definition}
\newtheorem{lemma}{Lemma}
\newtheorem{theorem}{Theorem}

% Blueprint macros (no-ops for LaTeX; metadata for isabelle-blueprint).
\newcommand{\isabelle}[1]{}
\newcommand{\isabelletheory}[1]{}
\newcommand{\isabellesession}[1]{}
\newcommand{\isabelleok}{}
\newcommand{\uses}[1]{}
\newcommand{\tags}[1]{}
\newcommand{\status}[1]{}
\newcommand{\blueprintstatus}[1]{}
\newcommand{\formalstatus}[1]{}
\newcommand{\agentstatus}[1]{}

\begin{document}

\title{Relations blueprint}
\maketitle

This blueprint develops the basic theory of binary relations and
equivalence relations. It demonstrates IsabelleBlueprint's LaTeX
ingestion path: each environment carries a \verb|\label|, an optional
\verb|\isabelle| fact, \verb|\uses| dependencies, and \verb|\tags|.
A \verb|\isabelleok| marker upgrades a named fact to ``found''.

\begin{definition}[Binary relation]
\label{relation}
\isabelle{Relations.relation_def}
\isabelleok
\tags{foundations}
A \emph{binary relation} on a set $A$ is a subset $R \subseteq A \times A$.
We write $a \mathrel{R} b$ for $(a, b) \in R$.
\end{definition}

\begin{definition}[Reflexive]
\label{reflexive}
\isabelle{Relations.reflexive_def}
\isabelleok
\uses{relation}
\tags{properties}
$R$ is \emph{reflexive} if $a \mathrel{R} a$ for every $a \in A$.
\end{definition}

\begin{definition}[Symmetric]
\label{symmetric}
\isabelle{Relations.symmetric_def}
\isabelleok
\uses{relation}
\tags{properties}
$R$ is \emph{symmetric} if $a \mathrel{R} b$ implies $b \mathrel{R} a$.
\end{definition}

\begin{definition}[Transitive]
\label{transitive}
\isabelle{Relations.transitive_def}
\uses{relation}
\tags{properties}
$R$ is \emph{transitive} if $a \mathrel{R} b$ and $b \mathrel{R} c$
imply $a \mathrel{R} c$.
\end{definition}

\begin{definition}[Equivalence relation]
\label{equivalence}
\isabelle{Relations.equivalence_def}
\uses{reflexive, symmetric, transitive}
\tags{foundations}
An \emph{equivalence relation} is a relation that is reflexive,
symmetric, and transitive.
\end{definition}

\begin{lemma}[Equality is an equivalence]
\label{eq-equivalence}
\isabelle{Relations.eq_equivalence}
\isabelleok
\uses{equivalence}
\tags{examples}
The equality relation $\{(a, a) \mid a \in A\}$ is an equivalence
relation on $A$.
\begin{proof}
Reflexivity, symmetry, and transitivity of $=$ are immediate.
\end{proof}
\end{lemma}

\begin{definition}[Equivalence class]
\label{eq-class}
\isabelle{Relations.eq_class_def}
\uses{equivalence}
\tags{quotients}
For an equivalence relation $R$ and $a \in A$, the \emph{equivalence
class} of $a$ is $[a]_R = \{ b \in A \mid a \mathrel{R} b \}$.
\end{definition}

\begin{theorem}[Classes partition the set]
\label{partition}
\isabelle{Relations.classes_partition}
\uses{eq-class, eq-equivalence}
\tags{quotients}
The equivalence classes of $R$ form a partition of $A$: every element
lies in exactly one class.
\begin{proof}
Reflexivity puts $a \in [a]_R$. If $[a]_R$ and $[b]_R$ share an element
then symmetry and transitivity force $[a]_R = [b]_R$.
\end{proof}
\end{theorem}

\end{document}
